\chapter{1975:John Cornforth and Vladimir Prelog}

\label{chap:chemistry1975}

The Nobel Prize in Chemistry for the year 1975 was awarded jointly to John Cornforth and Vladimir Prelog.

\textbf{Motivation:}

for his work on the stereochemistry of enzyme-catalyzed reactions (one half, John Cornforth) and for his research into the stereochemistry of organic molecules and reactions (the other half, Vladimir Prelog)

\section{John Cornforth}

\subsection*{Early Life and Education}

John Cornforth was born on 1917-09-07 in Sydney, Australia.

\subsection*{Career and Major Research}

Cornforth studied at the University of Sydney and at Oxford, where he received his doctorate. He worked at the National Institute for Medical Research in London and later became a professor at the University of Sussex. He was deaf from an early age. Using isotopic labeling, he worked out the stereochemistry of many enzyme-catalyzed reactions, particularly in the biosynthesis of cholesterol.

\subsection*{Nobel Prize-winning Work}

The work for which John Cornforth was awarded the Nobel Prize in Chemistry in 1975 was described by the Nobel Committee as: "for his work on the stereochemistry of enzyme-catalyzed reactions".

Cornforth traced the fate of individual hydrogen atoms and methyl groups through biosynthetic pathways, establishing the three-dimensional course of reactions catalyzed by enzymes.

\subsection*{Significance and Impact}

His methods for following stereochemistry with isotopic labels provided deep insight into how enzymes control the geometry of reactions.

\subsection*{Later Career and Honors}

Cornforth was knighted in 1977. He died on 2013-12-08 in Sussex, England.

\subsection*{Legacy}

Cornforth's stereochemical analyses are fundamental to mechanistic biochemistry and to the understanding of enzyme catalysis.

\subsection*{Selected Publications}

"Asymmetry and Enzyme Action" (Nobel Lecture, 1975)

\clearpage

\section{Vladimir Prelog}

\subsection*{Early Life and Education}

Vladimir Prelog was born on 1906-07-23 in Sarajevo, Bosnia.

\subsection*{Career and Major Research}

Prelog studied at the Czech Technical University in Prague, where he received his doctorate. He worked in the chemical industry in Zagreb and then joined the Swiss Federal Institute of Technology in Zurich, where he was professor of organic chemistry. He studied the stereochemistry of organic molecules, including medium-sized rings, and with Robert Cahn and Christopher Ingold developed the system for specifying the absolute configuration of chiral molecules.

\subsection*{Nobel Prize-winning Work}

The work for which Vladimir Prelog was awarded the Nobel Prize in Chemistry in 1975 was described by the Nobel Committee as: "for his research into the stereochemistry of organic molecules and reactions".

Prelog's research clarified the relationship between molecular structure and chirality, and the Cahn-Ingold-Prelog rules he helped devise are used to describe stereoisomers.

\subsection*{Significance and Impact}

The Cahn-Ingold-Prelog priority rules became the international standard for describing the configuration of stereocenters.

\subsection*{Later Career and Honors}

Prelog remained at the Swiss Federal Institute of Technology in Zurich until his retirement. He died on 1998-01-07 in Zurich, Switzerland.

\subsection*{Legacy}

The R/S system of stereochemical nomenclature developed with Prelog's participation is used throughout chemistry.

\subsection*{Selected Publications}

"Specification of Molecular Chirality" (Angewandte Chemie International Edition, 1966)

\clearpage

\section*{Chapter Summary}

The 1975 prize recognized Cornforth and Prelog for their work on the stereochemistry of enzyme-catalyzed reactions and of organic molecules. Their research established how molecular geometry governs chemical reactivity.

